% Impressum und Copyright-Seite (Dissertation)
% Bibliographische Angaben und Rechtsinformationen

% === Seitenstil ===
\thispagestyle{plain}

% === Prüfungskommission (bereits in Titelseite, hier als Referenz) ===
\vspace*{-1.0cm}
\begin{center}
  \textbf{Prüfungskommission:}
  
  \vspace{0.5cm}
  
  \begin{tabular}{rl}
    Vorsitzender: & Prof. Dr. Peter Hertling \\
    & \\
    1. Gutachter: & Prof. Dr. Gabi Dreo Rodosek\\
    2. Gutachter: & Prof. Dr. Oliver Rose\\
    & \\
    1. Prüfer: & Prof. Dr. Uwe Borghoff\\
    2. Prüfer: & Prof. Dr. Wolfgang Hommel\\
    3. Prüfer: & Prof. Dr.-Ing. Andreas Karcher\\
  \end{tabular}
\end{center}

\vfill

% === Bibliographische Informationen ===
\textbf{Bibliographische Informationen der Deutschen Bibliothek}

Die Deutsche Bibliothek verzeichnet diese Publikation in der Deutschen Nationalbibliographie; detaillierte bibliographische Daten sind im Internet über \url{http://dnb.ddb.de} abrufbar.

\vspace{0.5cm}

% === Copyright ===
\textcopyright{} 2017 Peter Hillmann

\vspace{0.5cm}

% === Herstellung und Verlag ===
Herstellung und Verlag: [Verlagsname einfügen]

Alle Rechte vorbehalten, insbesondere auch das Recht der fotomechanischen Wiedergabe, der Speicherung in elektronischen Medien, der Anfertigung sonstiger Kopien oder der Wiedergabe in anderen Medien wie im Internet. Ausnahmen sind nur mit schriftlicher Einwilligung des Autors möglich.

\vspace{0.5cm}

% === Haftungsausschluss ===
\textbf{Hinweis:}

Alle Angaben in diesem Buch wurden mit größter Sorgfalt erarbeitet und überprüft, trotzdem sind Fehler nicht auszuschließen. Daher kann weder der Herausgeber noch der Verlag eine Garantie oder Verantwortung für die Richtigkeit der Angaben übernehmen. Jegliche Haftung für Folgen, die auf die Verwendung von Angaben aus diesem Buch entstehen, muss ausgeschlossen werden. Für Hinweise auf Fehler sind Verlag und Herausgeber dankbar.

Die meisten im Werk genannten Produktbezeichnungen, Herstellernamen, Firmennamen und -logos sind marken-, warenzeichen- oder patentrechtlich geschützt. Die vorgestellten Verfahren und Programme werden ohne Rücksicht auf die Patentlage mitgeteilt. Ihre Angabe dient ausschließlich Lehrzwecken.

% === Open Source Hinweis (optional) ===
% Soweit im Buch Code von Programmen veröffentlicht wurde, die unter GNU General Public License (GNU GPL) stehen, richten sich die Nutzungsrechte an diesem Code ausschließlich nach der GNU GPL.

% === Template-Hinweise ===
% Für eigene Dissertationen anzupassen:
% - Copyright-Jahr und Name
% - Verlagsangaben
% - Prüfungskommission
% - Ggf. ISBN ergänzen